\section*{Capítulo \ref{ch:g10:soil-formation} --- O solo, um recurso de formação lenta}

\begin{solution}{exo:g10:soil-formation:1}
O \omterm{def:g10:soil-formation:chemical-weathering}{intemperismo químico} é a transformação dos minerais de uma rocha por
reações com a água e com os gases que ela leva (gás carbônico, oxigênio):
alguns minerais se dissolvem, outros se tornam minerais novos, estáveis na
superfície. A \omterm{def:g10:soil-formation:hydrolysis}{hidrólise} é o \omterm{def:g10:soil-formation:chemical-weathering}{intemperismo químico} de um silicato pela água e
pelos seus íons hidrogênio: os íons metálicos passam para a solução e o
restante do mineral é reconstruído num \omterm{def:g10:soil-formation:clay}{argilomineral}.
\end{solution}

\begin{solution}{exo:g10:soil-formation:2}
O \omterm{def:g10:soil-formation:weathering-profile}{saprolito} é rocha decomposta no lugar: mole e terroso, ele conserva os
grãos, as bandas e a estrutura da sua \omterm{def:g10:soil-formation:weathering-profile}{rocha-mãe} e passa para baixo, sem
ruptura, à rocha fresca. Um solo formado em areia de rio repousa sobre um
material trazido de outro lugar: os seus grãos são arredondados e
selecionados, e ele não mostra nada da estrutura da rocha que está por
baixo.
\end{solution}

\begin{solution}{exo:g10:soil-formation:3}
Os feldspatos se tornam argila (caulinita), liberando potássio, sódio ou
cálcio e sílica na água; a mica preta se torna argila e óxidos de ferro; o
quartzo fica intacto, como areia.
\end{solution}

\begin{solution}{exo:g10:soil-formation:4}
O ferro dos minerais escuros, \ce{Fe^2+}, é oxidado pelo oxigênio levado
pela água em óxidos de ferro(III) insolúveis: a goethita, amarelo-amarronzada,
e a hematita, vermelha, que recobrem os grãos e colorem todo o material.
\end{solution}

\begin{solution}{exo:g10:soil-formation:5}
\qty{0.017}{mm} por ano e \qty{1.537}{mm} por ano: $1.537/0.017 \approx 90$
vezes maior.
\end{solution}

\begin{solution}{exo:g10:soil-formation:6}
Cálcio: 1 e 1. Carbono: $1 + 1 = 2$ e 2. Oxigênio: $3 + 2 + 1 = 6$ e
$2 \times 3 = 6$. Hidrogênio: 2 e 2. Carga: 0, e $+2 - 2 = 0$.
\end{solution}

\begin{solution}{exo:g10:soil-formation:7}
A dissolução da calcita só produz íons dissolvidos, \ce{Ca^2+} e
\ce{HCO3-}: tudo é levado embora, e um calcário puro não deixa resíduo
sólido. A \omterm{def:g10:soil-formation:hydrolysis}{hidrólise} do feldspato produz um sólido, a caulinita, que fica; e
o quartzo do granito quase não reage. O granito deixa, portanto, argila e
areia.
\end{solution}

\begin{solution}{exo:g10:soil-formation:8}
As lâminas de argila e o húmus levam cargas negativas que retêm os íons
positivos de potássio \ce{K+} nas suas superfícies, de modo que a água da
chuva que passa não os leva embora. Uma raiz libera íons hidrogênio, que
tomam o lugar do potássio nas lâminas; o potássio liberado na água do solo é
absorvido pela raiz.
\end{solution}

\begin{solution}{exo:g10:soil-formation:9}
O \omterm{def:g10:soil-formation:chemical-weathering}{intemperismo químico} precisa de água, abundante nos climas úmidos e
escassa nos secos (e congelada nos frios); e as suas reações andam mais
depressa quando faz calor. Só com calor e água abundante, durante muito
tempo, as reações vão longe o bastante para remover quase tudo o que é
solúvel e deixar óxidos de ferro e de alumínio.
\end{solution}

\begin{solution}{exo:g10:soil-formation:10}
$10/0.017 \approx 590$ anos para \qty{1}{cm}; $300/0.017 \approx 17\,600$
anos para \qty{30}{cm}; $1000/0.017 \approx 59\,000$ anos para \qty{1}{m}.
\end{solution}

\begin{solution}{exo:g10:soil-formation:11}
Volume: $10\,000 \times 0.0015 = \qty{15}{m^3}$; massa: $15 \times 1200 =
\qty{18000}{kg}$, \num{18}~toneladas por hectare por ano.
\end{solution}

\begin{solution}{exo:g10:soil-formation:12}
$46\,000/0.017 \approx 2.7$ milhões de anos. Um \omterm{def:g10:soil-formation:weathering-profile}{saprolito} tão espesso exige
uma paisagem que permaneceu estável, sem que o seu manto intemperizado fosse
arrancado, durante milhões de anos.
\end{solution}

\begin{solution}{exo:g10:soil-formation:13}
Goethita: \ce{4 Fe^2+ + O2 + 6 H2O -> 4 FeO(OH) + 8 H+}: ferro 4 e 4,
oxigênio $2 + 6 = 8$ e $4 \times 2 = 8$, hidrogênio 12 e $4 + 8 = 12$, carga
$+8$ e $+8$. Hematita: \ce{4 Fe^2+ + O2 + 4 H2O -> 2 Fe2O3 + 8 H+}: ferro 4
e 4, oxigênio $2 + 4 = 6$ e 6, hidrogênio 8 e 8, carga $+8$ e $+8$.
\end{solution}

\begin{solution}{exo:g10:soil-formation:14}
Perda líquida $0.082 - 0.017 = \qty{0.065}{mm}$ por ano;
$300/0.065 \approx 4600$ anos. Sob a aração convencional, $300/1.52 \approx 200$ anos: mais de
vinte vezes menos.
\end{solution}

\begin{solution}{exo:g10:soil-formation:15}
A tonelada contém \qty{300}{kg} de ortoclásio. Cada quilograma dá
\qty{464}{g} de caulinita: $300 \times 0.464 \approx \qty{139}{kg}$.
Potássio: $300 \times 0.1405 \approx \qty{42}{kg}$.
\end{solution}

\begin{solution}{pb:g10:soil-formation:1}
\textbf{1.} $4 \times 0.3048 \approx \qty{1.22}{m}$.

\textbf{2.} $0.5 \times 30.48 \approx \qty{15}{cm}$.

\textbf{3.} O húmus, a matéria orgânica escura produzida pelos
decompositores a partir das plantas mortas.

\textbf{4.} Os óxidos de ferro, goethita e hematita, formados pela oxidação
do ferro(II) dos minerais do material de origem.

\textbf{5.} Não. Mais abaixo, o perfil passaria à rocha intemperizada, o
\omterm{def:g10:soil-formation:weathering-profile}{saprolito} (ou outro material intemperizado), e depois à \omterm{def:g10:soil-formation:weathering-profile}{rocha-mãe} fresca.

\textbf{6.} \ce{2 KAlSi3O8 + 2 H+ + 9 H2O -> Al2Si2O5(OH)4 + 2 K+ + 4 H4SiO4}.

\textbf{7.} $1000/278.33 \approx \qty{3.59}{mol}$.

\textbf{8.} $3.59/2 \approx \qty{1.80}{mol}$ de caulinita, $1.80 \times 258.16
\approx \qty{464}{g}$.

\textbf{9.} Os íons potássio e a sílica dissolvida: eles são levados pela
água aos rios, e pelos rios ao mar.

\textbf{10.} Ele quase não reage: os grãos de quartzo ficam como areia.

\textbf{11.} $1219/0.017 \approx 72\,000$ anos.

\textbf{12.} $1.537 - 0.017 = \qty{1.52}{mm}$ por ano.

\textbf{13.} $1219/1.52 \approx 800$ anos.

\textbf{14.} $72\,000/800 = 90$: o perfil é perdido noventa vezes mais
depressa do que se formou.

\textbf{15.} Não: $0.013 < 0.017$, a erosão é um pouco mais lenta que a
produção e o solo se mantém.

\textbf{16.} $T = S/(E - P)$.

\textbf{17.} $300/1.52 \approx 200$ anos.

\textbf{18.} $300/(0.082 - 0.017) = 300/0.065 \approx 4600$ anos.

\textbf{19.} Um solo perdido em alguns séculos leva dezenas de milhares de
anos para se formar de novo: na escala das vidas humanas e das sociedades, o
que se perde não é reposto.

\textbf{20.} $1000/1.52 \approx 660$ anos: \textbf{um metro de solo é
perdido em cerca de 660 anos, alguns séculos, sob a aração convencional}.
\end{solution}
